\documentclass[10pt,a4paper]{amsart}
\usepackage[margin=19mm]{geometry}
\usepackage{amsmath,amssymb,mathtools}
\usepackage[scr=rsfs,cal=euler]{mathalfa}
\usepackage{xfrac}
\usepackage[unicode,hidelinks,bookmarks=false]{hyperref}
\newcommand{\ie}{i.e.\ }
\newcommand{\cf}{cf.\ }
\newcommand{\resp}{respectively\ }
\newcommand{\Subsection}{\subsection}
\providecommand{\nobreakdash}{\nobreak\hbox{-}}
\setlength{\emergencystretch}{2em}
\multlinegap0pt
\usepackage{enumerate}
\usepackage{amssymb}
\usepackage{bm}
\usepackage{bbm}
\usepackage[shortlabels]{enumitem}
\usepackage{xcolor}
\usepackage{mathtools}
\newtheorem{lemma}{Lemma}
\newcommand{\Hil}{\mathcal{H}}
\newcommand{\X}{\mathfrak{X}}
\title{An explicit construction of heat kernels and Green's functions in measure spaces}
\author{Palle Jorgensen, Jay Jorgenson, Lejla Smajlovic}
\date{Source: arXiv:2512.24348v1 (CC BY 4.0)}
\begin{document}
\maketitle

\noindent\emph{Source and licence.} An explicit construction of heat kernels and Green's functions in measure spaces. Version arXiv:2512.24348v1 (CC BY 4.0), distributed by arXiv under the Creative Commons Attribution 4.0 International licence, http://creativecommons.org/licenses/by/4.0/, as recorded in the arXiv licence metadata for this exact version and shown on the abstract page https://arxiv.org/abs/2512.24348v1. Full licence notice: \texttt{licenses/arxiv-2512.24348-CC-BY-4.0.txt}.\par\smallskip

\noindent\textit{Editorial scope.} Counted unit: the article's lemma S (source0.tex lines 1537--1547). The article's complete original proof of it is delivered with it (source0.tex lines 1549--1622, one \texttt{proof} environment of 74 lines). No mathematics is written, repaired or re-derived: the statement and the proof are the article's own lines, in the article's own notation, and no retained line is substituted or re-flowed.  Retained uncounted as the source-local context the proof needs: lines 497--509.  Nothing was omitted from any retained span, so the package carries no omission record.  No retained line contains a citation, so the wrapper carries no bibliography and no bibliography entry.  No retained line contains a figure environment, an image-inclusion command or drawing code, so no image is delivered and the package declares no asset dependency.  Editorial formatting only, in addition: an AMS \texttt{amsart} preamble that restates the article's own declarations and macros used by the retained lines, namely the environments \texttt{lemma} on separate counters, together with the standard packages those lines need.  The retained environments print in this document as Lemma 1 in order of appearance, whereas the article prints the same items with the numbering of its own full document.  The complete native source of the article as distributed in the arXiv e-print for this exact version is bundled unchanged as \texttt{sources/191884/source0.tex}.
\begin{lemma} \label{lem: dif of inner product wrt t}
   Let $f:\X\times(0,\infty)$ be a function such that: (i) for any $t\in(0,\infty)$, function $x\mapsto f(x,t)$ belongs to $\Hil$, and (ii) for any $x\in\X$, function $t\mapsto f(x,t)$ is differentiable. Further assume that
   $$
   \partial_t f(x,t)\in \Hil
   \,\,\,\,\,
   \text{\rm for all $x\in\X$.}
   $$
    Then
   $$\partial_t \langle f(x,t),g(x)\rangle_{\Hil} = \langle \partial _t f(x,t),g(x)\rangle_{\Hil}
   \,\,\,\,\,
   \text{\rm for all $g\in\Hil.$}
   $$
\end{lemma}

\begin{lemma}
\label{lem: L of convolHilb S}
Let $H$ be an $\mathcal{H}$ parametrix for the heat operator on $(\X,\mu)$ of any order.
Let $f=f(x,y;t):\X\times \X\times\mathbb{R}_{> 0}\to\mathbb{R}$
be a continuous function in $t$ for all $x,y\in \X$. Assume further that for all $t\in \mathbb{R}_{>0}$ the function
$f(\cdot,y;t)$ when viewed as a function on $\X$, for all $y\in\X$ belongs to $\mathcal{H}$.  Then
$$
L_{x}( H\ast f)(x,y;t)=f(x,y;t)+(L_{x}H \ast f)(x,y;t)
$$
for all $x,y\in \X$ and $t\in\mathbb{R}_{>0}$.
\end{lemma}

\begin{proof} From the assumptions, we have that $f(\cdot, y;t)\in \mathcal{H}$ and $H(x,\cdot;t), \, \Delta_{x}H(x,\cdot;t)\in \mathcal{H}$
for all $(x,y;t) \in \X\times \X \times \mathbb{R}_{>0}$ are continuous in variable $t$.  Therefore,
their inner product is integrable on $[0,t]$, which proves that convolutions $H\ast f$ and $(\Delta_{x}H)\ast f$ are well defined.
Next, we claim that
\begin{equation}\label{eq. delta od conv}
  \Delta_{x}(H\ast f)(x,y;t) =  \left( (\Delta_{x}H)\ast f\right)(x,y;t).
\end{equation}
In order to prove \eqref{eq. delta od conv}, let us write
$$
H(x,\cdot;t)- H(z,\cdot;t)=\sum_{j\geq 1}h_{x,j}(z;t)\varphi_j(\cdot)
\,\,\,\,
\text{\rm and}
\,\,\,\,
f(\cdot,y;t)=  \sum_{j\geq 1}g_j(y;t)\varphi_j(\cdot)
$$
Then
\begin{align} \notag
 \Delta_{x}(H\ast f)(x,y;t) &= \int_\X\left(\int_0^t \left[ \langle H(x,\cdot;t)- H(z,\cdot;t-\tau),  f(\cdot,y;\tau) \rangle_{\mathcal{H}} \right] d\tau\right)\rho_x(dz)\\ \notag
 &=\int_\X\left(\int_0^t \left[ \sum_{j\geq 1} h_{x,j}(z;t-\tau)g_j(y;\tau) \right] d\tau\right)\rho_x(dz)\\ &= \int_0^t \left( \sum_{j\geq 1} \left(\int_\X  h_{x,j}(z;t-\tau) \rho_x(dz)\right) g_j(y;\tau)  \right)d\tau\label{eq. interchange od deltax}.
\end{align}
To be precise, the above computations involved an
interchange of the sum and the integral which is justified because
\begin{align*}
\sum_{j\geq 1} \int_\X \int_0^t &| h_{x,j}(z;t-\tau)g_j(y;\tau)  | \rho_x(dz) d\tau = \int_\X \int_0^t \left(\sum_{j\geq 1}| h_{x,j}(z;t-\tau)g_j(y;\tau)  |\right)  \rho_x(dz) d\tau\\ &\leq \int_\X \int_0^t \left(\sum_{j\geq 1} h_{x,j}(z;t-\tau)^2 \right)^{1/2} \left(\sum_{j\geq 1} g_{,j}(y;\tau)^2 \right)^{1/2 }\rho_x(dz) d\tau \\&=\int_\X \int_0^t ||H(x,\cdot;t-\tau)- H(z,\cdot;t-\tau) ||_\mathcal{H} ||f(\cdot,y;\tau) ||_\mathcal{H}\rho_x(dz) d\tau <\infty,
\end{align*}
where we used that the measure $\rho_x$ is finite and all norms are bounded.

Continuing, we claim that
\begin{equation}\label{eq. intermediate}
\int_\X  h_{x,j}(z;t-\tau) \rho_x(dz) = \langle \Delta_xH(x,\cdot;t-\tau),\varphi_j(\cdot) \rangle_\mathcal{H}.
\end{equation}
Since
\begin{equation}\label{eq.DeltaH_formula}
\Delta_xH(x,\cdot;t-\tau) = \int_\X \left(\sum_{\ell\geq 1} h_{x,\ell}(z;t-\tau) \varphi_\ell (\cdot)\right) \rho_x(dz),
\end{equation}
in order to prove \eqref{eq. intermediate} it suffices to show that one can interchange the sum and the integral in \eqref{eq.DeltaH_formula}.
This easily follows from the Cauchy-Schwarz inequality, namely
$$
\sum_{\ell\geq 1} \int_\X \left| h_{x,\ell}(z;t-\tau) \varphi_\ell (\cdot) \right| \rho_x(dz)\leq \int_\X \left( \sum_{\ell\geq 1} h_{x,\ell}(z;t-\tau)^2\right)^{1/2}  \left( \sum_{\ell\geq 1} \varphi_\ell(\cdot)^2\right)^{1/2} <\infty.
$$

Now, when combining \eqref{eq. interchange od deltax} and \eqref{eq. intermediate} we arrive at
\begin{align*}
\Delta_{x}(H\ast f)(x,y;t) &= \int_0^t \left( \sum_{j\geq 1} \langle \Delta_xH(x,\cdot;t-\tau),\varphi_j(\cdot) \rangle_\mathcal{H}  g_j(y;\tau)  \right)d\tau \\&=  \int_0^t \left( \langle \Delta_x H(x,\cdot;t-\tau),f(\cdot,y;\tau) \rangle_\mathcal{H}  \right)d\tau \\&=  \left( (\Delta_{x}H)\ast f\right)(x,y;t).
\end{align*}
Therefore,
\begin{align}\notag
L_{x}(H\ast f)(x,y;t)&=
    \frac{\partial}{\partial t}(H\ast f)(x,y;t)
    + \Delta_{x}(H\ast f)(x,y;t)\\ &=
  \frac{\partial}{\partial t}(H\ast f)(x,y;t)+  \left( (\Delta_{x}H)\ast f\right)(x,y;t).\label{eq:heat_of_convolution2}
\end{align}
The function $\langle H(x,z;t-r),f(z,y;r) \rangle_{\mathcal{H}}$ is continuous in the time variable.
This follows because the integrand is continue due to the H\"older inequality and the assumptions on the parametrix, namely boundedness
uniform boundedness in $r\in[0,t]$ by an integrable function.  So we can apply the
Leibniz integration formula.  Upon doing so, we obtain that
the first term on the right hand side of \eqref{eq:heat_of_convolution2} is equal to
\begin{align} \label{eq.deriv in t}
\lim_{\varepsilon\to 0^+}\frac{\partial}{\partial t}&\int\limits _{0}^{t-\varepsilon}\left( \langle H(x,\cdot;t-\tau),f(\cdot,y;\tau) \rangle_\mathcal{H}  \right)d\tau \\&
=\lim_{\varepsilon\to 0^+}\left( \langle H(x,\cdot;\varepsilon),f(\cdot,y;\tau) \rangle_\mathcal{H}  \right) +\int\limits _{0}^{t} \frac{\partial}{\partial t}\left( \langle H(x,\cdot;t-\tau),f(\cdot,y;\tau) \rangle_\mathcal{H}  \right)d\tau .\notag
\end{align}
In view of assumption (4) on the parametrix $H$, we get from \eqref{eq.deriv in t}, combined with Lemma \ref{lem: dif of inner product wrt t} that
$$
\frac{\partial}{\partial t}\int\limits _{0}^{t}\int\limits_{\X}H(x,z;t-r)f(z,y;r) \mu(dz)dr= f(x,y;t) + \left(\frac{\partial}{\partial t}H\ast f\right)(x,y;t).
$$
Therefore,
\begin{align*}
L_{x}(H\ast f)(x,y;t) &=
  \frac{\partial}{\partial t}(H\ast f)(x,y;t)+   (\Delta_{x}H\ast f)(x,y;t)
\\&=f(x,y;t)+\left(\frac{\partial}{\partial t}H\ast f\right)(x,y;t)+(\Delta_{x}H\ast f)(x,y;t)
\\&=f(x,y;t)+(L_{x}H\ast f)(x,y;t),
\end{align*}
as claimed.
\end{proof}

\end{document}
